\documentclass[10pt]{scrartcl}
\usepackage[a4paper, total={5.5in, 8in}]{geometry}
\usepackage[usenames,dvipsnames,svgnames]{xcolor}
\usepackage[shortlabels]{enumitem}
\usepackage[framemethod=TikZ]{mdframed}
\usepackage{amsmath,amssymb,amsthm}
\usepackage[colorlinks]{hyperref}

\hypersetup{
	colorlinks=true,
	linkcolor=blue,
	filecolor=blue,      
	urlcolor=blue,
	pdftitle={MAT 362}
}

\usepackage{microtype}
\usepackage{mathtools}
\usepackage[headsepline]{scrlayer-scrpage}
\usepackage{thmtools}
\usepackage[textsize=small]{todonotes}

\addtolength{\textheight}{3.14cm}
\providecommand{\ol}{\overline}
\providecommand{\eps}{\varepsilon}
\providecommand{\half}{\frac{1}{2}}
\providecommand{\dang}{\measuredangle} %% Directed angle
\providecommand{\CC}{\mathbb C}
\providecommand{\EE}{\mathbb E}
\providecommand{\FF}{\mathbb F}
\providecommand{\NN}{\mathbb N}
\providecommand{\QQ}{\mathbb Q}
\providecommand{\RR}{\mathbb R}
\providecommand{\ZZ}{\mathbb Z}
\providecommand{\dg}{^\circ}
\providecommand{\ind}[1]{\mathbf 1_{#1}}
\providecommand{\ii}{\item}
\providecommand{\alert}[1]{{\sffamily\textbf{\textcolor{magenta}{#1}}}}
\providecommand{\opname}{\operatorname}
\providecommand{\li}{\opname{li}}
\providecommand{\ls}{\opname{ls}}
\providecommand{\inv}{^{-1}}
\providecommand{\setdiff}{\smallsetminus}
\providecommand{\mb}[1]{\mathbf{#1}}
\providecommand{\abs}[1]{\left|{#1}\right|}
\providecommand{\symdiff}{\Delta}
\providecommand{\jim}{m_{*,(J)}}
\providecommand{\jom}{m^{*,(J)}}
\providecommand{\pcint}{\mathrm{p.c.}\int}
\providecommand{\dif}{\;d}
\providecommand{\dist}{\operatorname{dist}}
\providecommand{\diam}{\operatorname{diam}}
\renewcommand{\hat}{\widehat}
\providecommand{\norm}[1]{\| #1 \|}

\reversemarginpar
\setlength{\marginparwidth}{3cm}

\mdfdefinestyle{mdbluebox}{roundcorner=10pt,innerbottommargin=9pt,
	linecolor=blue,backgroundcolor=TealBlue!5,}
\declaretheoremstyle[headfont=\sffamily\bfseries\color{MidnightBlue},
mdframed={style=mdbluebox},]{thmbluebox}

\mdfdefinestyle{mdredbox}{frametitlefont=\bfseries,innerbottommargin=8pt,
	nobreak=true,backgroundcolor=Salmon!5,linecolor=RawSienna,}
\declaretheoremstyle[headfont=\bfseries\color{RawSienna},
mdframed={style=mdredbox},headpunct={\\[3pt]},postheadspace=0pt,]{thmredbox}

\mdfdefinestyle{mdgreenbox}{linecolor=ForestGreen,backgroundcolor=ForestGreen!5,
	linewidth=2pt,rightline=false,leftline=true,topline=false,bottomline=false,}
\declaretheoremstyle[headfont=\bfseries\sffamily\color{ForestGreen!70!black},
mdframed={style=mdgreenbox},headpunct={ --- },]{thmgreenbox}

\mdfdefinestyle{mdorangebox}{linecolor=RedOrange,backgroundcolor=RedOrange!5,
	linewidth=2pt,rightline=false,leftline=true,topline=false,bottomline=false,}
\declaretheoremstyle[headfont=\bfseries\sffamily\color{RedOrange!70!black},
mdframed={style=mdorangebox},headpunct={ --- },]{thmorangebox}

\mdfdefinestyle{mdblackbox}{linecolor=black,backgroundcolor=RedViolet!5!gray!5,
	linewidth=3pt,rightline=false,leftline=true,topline=false,bottomline=false,}
\declaretheoremstyle[mdframed={style=mdblackbox}]{thmblackbox}


\declaretheorem[style=thmredbox,name=Problem,numberwithin=section]{problem}

\declaretheorem[style=thmredbox,name=Example,sibling=problem]{example}
\declaretheorem[style=thmredbox,name=$\bigstar$ Problem,sibling=problem]{niceprob}

\declaretheorem[style=thmbluebox,name=Theorem,sibling=problem]{theorem}
\declaretheorem[style=thmbluebox,name=Corollary,sibling=theorem]{corollary}
\declaretheorem[style=thmbluebox,name=Proposition,sibling=theorem]{prop}
\declaretheorem[style=thmbluebox,name=Lemma,sibling=theorem]{lemma}
\declaretheorem[style=thmbluebox,name=Theorem,numbered=no]{theorem*}
\declaretheorem[style=thmbluebox,name=Lemma,numbered=no]{lemma*}
\declaretheorem[style=thmgreenbox,name=Claim,numberwithin=theorem]{claim}
\declaretheorem[style=thmgreenbox,name=Claim,numbered=no]{claim*}
\declaretheorem[style=thmblackbox,name=Remark,sibling=theorem]{remark}
\declaretheorem[style=thmblackbox,name=Remark,numbered=no]{remark*}
\declaretheorem[style=thmgreenbox,name=Definition,sibling=theorem]{definition}
\declaretheorem[style=thmgreenbox,name=Definition,numbered=no]{definition*}
\declaretheorem[style=thmorangebox,name=Warning,sibling=theorem]{warning}
\declaretheorem[style=thmorangebox,name=Warning,numbered=no]{warning*}
\declaretheorem[style=thmorangebox,name=Note,numbered=no]{note}
\declaretheorem[style=thmblackbox,name=Exercise,sibling=problem]{exercise}
\declaretheorem[style=thmblackbox,name=Exercise,numbered=no]{exercise*}

\newenvironment{soln}{\begin{proof}[Solution]}{\end{proof}}
\newenvironment{walkthrough}{\noindent \textbf{Walkthrough.}}{\medskip}
\newlist{walk}{enumerate}{3}
\setlist[walk]{label=\bfseries (\alph*)}

\providecommand{\pow}{\operatorname{Pow}}
\providecommand{\vocab}[1]{{\textbf{\textcolor{ForestGreen}{#1}}}}
\providecommand{\lcm}{\operatorname{lcm}}
\providecommand{\ord}{\operatorname{ord}}
\providecommand{\ex}{\opname{ex}}
\providecommand{\floor}[1]{\left\lfloor #1\right\rfloor}
\providecommand{\ceil}[1]{\left\lceil #1\right\rceil}

\usepackage{asymptote}
\begin{asydef}
	defaultpen(fontsize(10pt));
	unitsize(1cm);
	size(8cm); // set a reasonable default
	usepackage("amsmath");
	usepackage("amssymb");
	settings.tex="pdflatex";
	settings.outformat="pdf";
	// Replacement for olympiad+cse5 which is not standard
	import geometry;
	// recalibrate fill and filldraw for conics
	void filldraw(picture pic = currentpicture, conic g, pen fillpen=defaultpen, pen drawpen=defaultpen)
	{ filldraw(pic, (path) g, fillpen, drawpen); }
	void fill(picture pic = currentpicture, conic g, pen p=defaultpen)
	{ filldraw(pic, (path) g, p); }
	// some geometry
	pair foot(pair P, pair A, pair B) { return foot(triangle(A,B,P).VC); }
	pair orthocenter(pair A, pair B, pair C) { return orthocentercenter(A,B,C); }
	pair centroid(pair A, pair B, pair C) { return (A+B+C)/3; }
	// cse5 abbreviations
	path CP(pair P, pair A) { return circle(P, abs(A-P)); }
	path CR(pair P, real r) { return circle(P, r); }
	pair IP(path p, path q) { return intersectionpoints(p,q)[0]; }
	pair OP(path p, path q) { return intersectionpoints(p,q)[1]; }
	path Line(pair A, pair B, real a=0.6, real b=a) { return (a*(A-B)+A)--(b*(B-A)+B); }
	// cse5 more useful functions
	picture CC() {
		picture p=rotate(0)*currentpicture;
		currentpicture.erase();
		return p;
	}
	pair MP(Label s, pair A, pair B = plain.S, pen p = defaultpen) {
		Label L = s;
		L.s = "$"+s.s+"$";
		label(L, A, B, p);
		return A;
	}
	pair Drawing(Label s = "", pair A, pair B = plain.S, pen p = defaultpen) {
		dot(MP(s, A, B, p), p);
		return A;
	}
	path Drawing(path g, pen p = defaultpen, arrowbar ar = None) {
		draw(g, p, ar);
		return g;
	}
\end{asydef}

\ihead{\footnotesize\textbf{MAT 362}}
\ohead{\footnotesize \today}

\title{\vspace{-2.0cm} MAT 362 Spring 2024 Notes}
\author{\large Aditya Pahuja}
\date{\large Test on 3/11!}
\begin{document}
\maketitle
\tableofcontents
\pagebreak

\section{Curves}
\renewcommand{\thesubsection}{\S}
\subsection{1.2, 1.3 Parametrized curves and regular curves}
We start with rehashing a some of the ideas
from basic multivariable calculus.

\begin{definition}[Parametrized differentiable curve]
	A \vocab{parametrized differentiable curve}
	is a differentiable map $\alpha\colon I\to\RR^3$
	from an open interval $(a,b)$ to $\RR^3$.
	Note that we will be using ``differentiable'' to mean that
	at all points in $I$, $\alpha$ has partials of all orders.
\end{definition}

We choose to think of a curve with respect to its parametrization
rather than just a subset of $\RR^3$; the set $\alpha(I)$ is the
\vocab{trace} of $\alpha$.
The derivative of a differentiable curve $\alpha$ at $t$
is the tangent vector/velocity vector at $t$.

\begin{exercise}
	If $u(t)$, $v(t)$ are parametrized differentiable curves on $I$,
	show that $(u\cdot v)' = u\cdot v' + u'\cdot v$, where $\cdot$
	is the dot product.
\end{exercise}

\begin{definition}[Regular curve]
	A parametrized differentiable curve $\alpha\colon I\to\RR^3$
	is \vocab{regular} if $\alpha'(t)\ne 0$ for all $t\in I$.
	If $\alpha'(t) = 0$ then we call $t$ a \vocab{singular point}
	of $\alpha$.
\end{definition}

The arc length of a regular curve from times $t_0$ to $t$ is 
\[s(t) := \int_{t_0}^t |\alpha'(x)|\dif x.\]
By the fundamental theorem of calculus, $s$ is differentiable
and $s'(t) = |\alpha'(t)|$
(note that we require $\alpha'(t)\ne 0$ to conclude differentiability
of all orders here because absolute value isn't differentiable at zero).

In the event that $s(t) = t$ for all $t$ (i.e. the parameter $t$
is equal to the arc length measured off of $t_0$), we thus have
$ds/dt = 1 = |\alpha'(t)|$, so the tangent vector has constant velocity.
Conversely if $|\alpha'(t)| = 1$ then
\[s = \int_{t_0}^t \dif x = t - t_0.\]
That is, $|\alpha'(t)| = 1$ for all $t$ is equivalent to $t$
measuring the arc length off a given point.
We are very interested in this case of the curve
being parametrized by arc length, since it makes our computations much nicer.

\subsection{1.4 Vector product}
The vector product is just the cross product:
\begin{definition}
	The \vocab{vector product} $u\wedge v$ of vectors $u,v\in\RR^3$ is
	an anticommutative, nonassociative bilinear operation on $\RR^3\times\RR^3$
	defined as the unique vector such that
	\[(u\wedge v)\cdot w = \det(u, v, w)\]
	for all $w\in\RR$.
\end{definition}

\begin{exercise}[Hope you remember how cross product works]
	(i) When is $u\wedge v$ zero?
	
	(ii) Given $u$ and $v$ with $u\wedge v\ne 0$,
	describe the set of all $w$ such that $(u\wedge v)\cdot w = 0$.
	
	(iii) Give an example to show that the vector product
	is not associative.
\end{exercise}

One can show that
\[(u\wedge v)\cdot (x\wedge y) = \begin{vmatrix}
	u\cdot x&v\cdot x\\ u\cdot y&v\cdot y
\end{vmatrix},\]
so that $|u\wedge v|^2 = (u\wedge v)\cdot(u\wedge v)$ is the area
of the parrallelogram spanned by $u$ and $v$ (check this ---
recall that $|u\cdot v| = |u||v||\cos\theta|$ where $u,v$ make angle $\theta$).

There is a nice differentiation rule for the vector product,
in analogy to the product rule for the dot product (check it!):
\[(u\wedge v)' = u\wedge v' + u'\wedge v.\]

\subsection{1.5 Curvature and shit}
Throughout this section let $\alpha\colon I\to\RR^3$ be a regular curve
parametrized by arc length.

\begin{definition}[Curvature]
	The \vocab{curvature} of $\alpha$ is $\kappa(s) := |\alpha''(s)|$.
\end{definition}

Since $\alpha'$ has constant length, one can show (exercise!)
that $\alpha''$ is normal to $\alpha'$, so $\kappa$ is a measure of
how quickly the curve is pulling away from its tangent vector,
i.e. how sharply the curve is turning at $s$.

\begin{exercise}
	Show that $\alpha$ has curvature zero everywhere
	if and only if it is a straight line.
\end{exercise}

The \vocab{normal vector} $n(s) = \alpha''(s)/\kappa(s)$ to $\alpha$
at $s$ is a unit vector normal to $t(s) := \alpha'(s)$
(so $t$ is the unit tangent vector at $s$).
The plane determined by $\alpha'$ and $n$ at a point $s$
is the \vocab{osculating plane} of $\alpha$ at $s$;
intuitively, this is the plane that $\alpha$ ``most closely'' hugs
as $s$ varies.
Note that the osculating plane doesn't exist if the curvature at $s$ is zero
(intuitively, there are infinitely many candidates for the osculating plane).
From this point on we shall restrict ourselves to assuming that
$\alpha''$ is everywhere nonzero (i.e. there are no
``singular points of order $1$'').

The unit vector $b(s) = t(s)\wedge n(s)$ is the \vocab{binormal vector}
at $s$; it is a unit vector normal to the osculating plane.
The quantity $|b'(s)|$, then, measures how quickly
the curve pulls away from its osculating plane at $s$.
Note that $b' = t'\wedge n + t\wedge n' = t\wedge n'$ is normal to $t$ because
because $t' = \kappa n$ and that $b'$ is normal to $b$ because $|b|$ is fixed.
Therefore $b'$ is parallel to $n$; we write $b'(s) = \tau(s)n(s)$,
where $\tau(s)$ is called the \vocab{torsion} at $s$.

\begin{exercise}
	Under the assumption $\kappa\ne 0$ everywhere, show that
	$\alpha(I)$ is contained in a plane if and only if the torsion
	is zero everywhere.
\end{exercise}

To sum up, we have three pairwise orthogonal unit vectors $t$, $n$, $b$,
which form the so-called ``Frenet trihedron,'' which
satisfy the ``Frenet formulas''
\begin{align*}
	t' &= \kappa n\\
	n' &= -\kappa t-\tau b\\
	b' &= \tau n
\end{align*}
where the second equation comes from differentiating $n=b\wedge t$.
A small amount of terminology:
The $tn$-plane is the osculating plane,
the $nb$-plane is the \vocab{normal plane}, and
the $tb$-plane is the \vocab{rectifying plane}.
Also, the \vocab{radius of curvature} at $s$ is defined as $1/\kappa$,
which one can think of intuitively as
the radius of the circle that ``best approximates'' the curve near $s$.

The extremely nice thing about all this theory is the following fact,
which roughly says that curvature and torsion uniquely
determine a curve up to rigid motion:
\begin{theorem}[Fundamental theorem of the local theory of curves]
	Given differentiable functions $\kappa(s)>0$ and $\tau(s)$,
	where $s$ varies over an interval $I$, there is a
	regular parametrized curve $\alpha\colon I\to\RR^3$ such that
	$s$ is the arc length, $\kappa(s)$ is its curvature,
	and $\tau(s)$ is its torsion.
	Moreover, if $\ol\alpha$ is another such curve, then
	there exists an orthogonal linear map $\rho$ on $\RR^3$
	with positive determinant and a vector $c$
	such that $\ol\alpha=\rho\circ \alpha + c$.
\end{theorem}

\subsection{1.7 Isoperimetric inequality}
I will not reproduce the proof of the isoperimetric inequality here.

A \vocab{closed plane curve} is a parametrized differentiable curve
$\alpha\colon[a,b]\to\RR^2$ such that $\alpha(a)=\alpha(b)$
and $\alpha^{(n)}(a) = \alpha^{(n)}(b)$ for all integers $n$,
i.e. the derivatives of all orders agree at $a$ and $b$.
A closed plane curve is \vocab{simple} if it is injective on $[a,b)$,
so its only self intersection is at the endpoints.
The Jordan curve theorem says that every simple closed curve
partitions the plane into a bounded ``inside'' region
and an unbounded ``outside'' region.
\begin{theorem}[Isoperimetric inequality]
	Let $C$ be a simple closed curve with length $\ell$
	and let $A$ be the area of the region bounded by $C$. Then
	\[\ell^2-4\pi A\ge 0,\]
	with equality if and only if $C$ is a circle.
\end{theorem}

\section{Regular surfaces}
\subsection{2.2 What's a regular surface?}
We want to now develop some tools that help us understand
the properties of surfaces. First, we need to do some work
to more rigorously describe what a surface is,
or at least to define a class of ``sufficiently nice'' surfaces.
One way to look at this is to consider sets that locally look like $\RR^2$,
in the following sense:
\begin{definition}[Regular surface]
	A subset $S\subseteq\RR^3$ is a \vocab{regular surface}
	if, for each $p\in S$, there exists a neighborhood
	$V\subseteq \RR^3$ of $p$ and an open set $U\subseteq\RR^2$
	for which there is a map $\mb x\colon U\to V\cap S$ satisfying
	\begin{enumerate}[(i)]
		\ii $\mb x$ is differentiable, meaning $\mb x$
		has partials of all orders.
		\ii $\mb x$ is a homeomorphism between $U$ and $V\cap S$,
		meaning $\mb x$ has a continuous inverse; it is continuous already
		due to condition (i).
		\ii For each $q\in U$, the differential $d\mb x_q\colon\RR^2\to\RR^3$
		is injective.
	\end{enumerate}
	We say that $\mb x$ is a \vocab{parametrization} of $S$
	in a neighborhood of $p$, and we call the neighborhood
	$V\cap S$ of $p$ in $S$ a \vocab{coordinate neighborhood}.
\end{definition}

We will show later that this is more naturally interpreted as
saying that $S$ is ``locally diffeomorphic'' to (an open subset of) $\RR^2$.
The third condition can be reinterpreted as saying that
the columns of the Jacobian matrix of $\mb x(u, v)$ are independent;
recall that the Jacobian matrix is
\[d\mb x_q = \begin{pmatrix}
	\partial x/\partial u(q) & \partial x/\partial v(q)\\
	\partial y/\partial u(q)&\partial y/\partial v(q)\\
	\partial z/\partial u(q)&\partial z/\partial v(q)
\end{pmatrix}\]

\begin{example}[Sphere is regular surface]
	The unit sphere $S^2$ is a regular surface. There are many ways to exhibit
	a covering of coordinate neighborhoods of a sphere;
	perhaps the most natural way is to note that the intersection of
	$S^2$ with each of the half-spaces determined by one of
	the coordinate planes (e.g. $S^2\cap \{(x,y,z)\in\RR^3 : z>0\}$)
	is a coordinate neighborhood; for example,
	$S^2\cap \{(x,y,z)\in\RR^3 : z>0\}$ admits a parametrization as
	\[\mb x(u, v) = (u, v,\sqrt{1-u^2-v^2})\]
	while $S^2\cap \{(x,y,z)\in\RR^3 : y<0\}$ can be parametrized as
	\[\mb x(u, v) = (u, -\sqrt{1-u^2-v^2},v).\]
	It is not too hard to check that these satisfy
	the conditions of the definition of a regular surface
	and that the six coordinate neighborhoods shown
	cover the entire sphere.
\end{example}

\begin{exercise}
	Check that a cylinder, the locus of points satisfying $x^2 + y^2 = 1$,
	is a regular surface.
\end{exercise}

The parametrizations of the coordinate neighborhoods in the example of $S^2$
generalize in an obvious way:
\begin{prop}
	If $f\colon U\to \RR$ is differentiable on an open subset of $\RR^2$,
	then the set $\{(x,y,f(x,y))\in\RR^3 : (x,y)\in U\}$
	is a regular surface.
\end{prop}

\begin{exercise}
	Prove this.
\end{exercise}

However, there is some sense in which the sphere or the cylinder
are ``obviously'' regular that makes the amount of work necessary
to check their regularity unsatisfying.
It turns out that there is a nice condition for when
the set of points satisfying $F(x, y, z) = c$ for some
differentiable $F$ and constant $c$ forms a regular surface.

\begin{definition}
	Let $F\colon U\subseteq\RR^n\to\RR^m$ be a differentiable map
	defined on an open subset $U$ of $\RR^n$.
	We say that $p\in U$ is a \vocab{critical point} of $F$
	if $dF_p\colon\RR^n\to\RR^m$ is not surjective.
	The image $F(p)\in\RR^m$ is a \vocab{critical value} of $F$.
	A point of $\RR^m$ which is not a critical value
	is called a \vocab{regular value} of $F$.
\end{definition}

\begin{prop}[Regular values generate regular surfaces]
	If $f\colon U\subseteq\RR^3\to\RR$ is differentiable
	and $a\in f(U)$ is a regular value of $f$, then $f\inv(a)$
	is a regular surface.
\end{prop}

Note that $a$ is a regular value of $f$ if and only if
$f_x$, $f_y$, and $f_z$ are not all simultaneously zero
at some point in the set $f\inv(a)$.

\begin{proof}[Proof]
	For some point $p=(x_0,y_0,z_0)\in f\inv(a)$,
	suppose WLOG that $f_z(p)\ne 0$.
	Then, by the implicit function theorem, there is a neighborhood
	$U\subseteq\RR^2$ of $(x_0,y_0)$ and a neighborhood $V\subseteq\RR$ of $z_0$
	for which there is a unique differentiable $g\colon U\to V$
	such that $g(x_0,y_0) = z_0$ and $f(x, y, g(x,y)) = a$
	for all $(x,y)\in U$.
	
	Now the map $\mb x\colon U\to\RR^3$
	given by $\mb x(u,v) = (u,v,g(u,v))$ is a parametrization
	of a neighborhood of $p$ in $f\inv(a)$
	because it is the graph of a differentiable function $g$
	on the open set $U$, so we are done
	for we have found a coordinate neighborhood $\mb x(U)\subseteq f\inv(a)$
	of each $p\in f\inv(a)$.
\end{proof}

(The book does a different proof with the inverse function theorem,
so maybe also read that if you want.)

\begin{example}
	The cylinder and sphere are now very clearly regular surfaces,
	as is the ellipsoid
	\[\frac{x^2}{a^2} + \frac{y^2}{b^2} + \frac{z^2}{c^2} = 1\]
	and the hyperboloid of two sheets
	\[-x^2-y^2+z^2=1\]
	and the torus
	\[z^2 = r^2 - (\sqrt{x^2+y^2} - a)^2\]
	for $a > r > 0$
	(exercise: show that this actually defines a torus!).
\end{example}

The following is a nice shortcut for determining when a map $U\to S$
is a parametrization if we already know a priori that $S$ is regular.
\begin{prop}
	Let $p\in S$ be a point of a regular surface $S$.
	Suppose that $\mb x\colon U\subseteq\RR^2\to\RR^3$
	is an injective differentiable map with injective differential
	such that $p\in\mb x(U)\subseteq S$.
	Then $\mb x\inv\colon \mb x(U)\to U$ is continuous.
\end{prop}

\begin{exercise}
	Prove this; I do not feel like writing down a proof rn
\end{exercise}

\subsection{2.3 Differentiable functions}
The point of all of the data carried by a regular surface
is that it very strongly respects a lot of the structure
of differential calculus, which is key if we are going to
try working with differential calculus in a more general context.

A differentiable function $f\colon S\to\RR$
for some regular surface $S$ can be defined very naturally:
\begin{definition}
	Let $S$ be a regular surface. Then $f\colon S\to\RR$ is differentiable
	at a point $p$ if, given a parametrization
	$\mb x\colon U\subseteq\RR^2\to\RR^3$ of a neighborhood of $p$,
	the function $f\circ\mb x\colon U\to\RR$ is differentiable
	at $\mb x\inv(p)$.
\end{definition}

There is an important technical note that we have to
resolve immediately, though: we need to ensure that
differentiability is independent of the choice of $\mb x$.

\begin{prop}[Change of parameters]
	Let $p$ be a point of a regular surface $S$ and let
	$\mb x\colon U\subseteq\RR^2\to S$,
	$\mb y\colon V\subseteq\RR^2\to S$
	be two parametrizations of $S$ such that $p\in\mb x(U)\cap\mb y(V)=W$.
	Then, the ``change of coordinates''
	$h = \mb x\inv\circ y\colon \mb y\inv(W)\to\mb x\inv(W)$
	is a diffeomorphism.
\end{prop}

\begin{exercise}
	Prove this; I'm lazy
	(if I had to guess, ya gotta use implicit function theorem.)
\end{exercise}

We can also talk about maps between surfaces being differentiable:
\begin{definition}[Differentiable map between surfaces]
	Let $S_1$ and $S_2$ be regular surfaces.
	Then, a continuous map $\varphi\colon V_1\subseteq S_1\to S_2$
	on an open subset $V_1$ of $S_1$ is differentiable at $p\in V_1$
	if, given parametrizations $\mb x_1\colon U_1\subseteq\RR^2\to S_1$
	and $\mb x_2\colon U_2\subseteq\RR^2\to S_2$ with $p\in\mb x_1(U_1)$ and
	$\varphi(p)\in\mb x_2(U_2)$, the map
	\[\mb x_2\inv\circ\varphi\circ\mb x_1\colon U_1\to U_2\]
	is differentiable at $\mb x_1\inv(p)$.
\end{definition}

The change of parameters proposition lets one easily show
that this definition is independent of the choice of parametrizations.

\begin{exercise}
	Let $S_1$ and $S_2$ be regular surfaces
	and assume that $S_1\subseteq V\subseteq\RR^3$
	where $V$ is open in $\RR^3$.
	Suppose that $\varphi\colon V\to\RR^3$ is differentiable
	and that $\varphi(S_1)\subseteq S_2$.
	Then, the restriction $\varphi|_{S_1}\colon S_1\to S_2$
	is a differentiable map.
\end{exercise}

Now, we can say that two surfaces $S_1$ and $S_2$ are
diffeomorphic if there is a differentiable function $S_1\to S_2$
with differentiable inverse.
One can show easily that $U$ and $\mb x(U)$ are diffeomorphic
for every parametrization $\mb x$ of a coordinate neighborhood
of a regular surface $S$. This is the true geometric intuition
behind what a regular surface looks like --- locally, its properties
relating to differentiability are the same as those of $\RR^2$.
Indeed, we could replace the definition of regular surface
with the more abstract ``a set that is locally diffeomorphic to $\RR^2$''
and recover all of this theory.

\subsection{2.3 Regular curves}
We can define regular curves in analogy to regular surfaces,
as subsets of $\RR^3$ that can be covered with coordinate neighborhoods
from $\RR$.

We can sometimes generate regular surfaces by displacing regular curves
in a certain way. For example, if $C$ is a regular connected plane curve
lying in the $xz$-plane, then parametrize $C$ as
$x=f(v)$, $z=g(v)$ for $a<v<b$, $f(v)>0$. Then the surface of revolution $S$
\[\mb x(u,v) = (f(v)\cos u, f(v)\sin u, u)\]
obtained by rotating the generating curve $C$ around the $z$-axis can be shown
to be a regular surface (exercise, check this).
The circles traced by fixed points of $C$ are \vocab{parallels} of $S$
(i.e. fixing $v$), while the images of $C$ under the rotations
are the \vocab{meridians} of $C$ (i.e. fixing $u$).
\pagebreak
\subsection{2.3 Regular parametrized surfaces}
We can instead define regular parametrized surfaces
in analogy with regular parametrized curves.
\begin{definition}
	A \vocab{parametrized surface} $\mb x\colon U\subseteq\RR^2\to\RR^3$
	is a differentiable map $\mb x$ on an open set $U$.
	We say that $\mb x$ is \vocab{regular} if $d\mb x_q\colon\RR^2\to\RR^3$
	is injective for all $q\in U$.
	A point where $d\mb x_q$ is not injective is a
	\vocab{singular point} of $\mb x$.
\end{definition}

\begin{example}
	Let $\alpha\colon I\to\RR^3$ be a nonplanar regular parametrized curve
	whose curvature is everywhere nonzero. Then we define the
	\vocab{tangent surface} of $\alpha$ to be the parametrized surface
	$\mb x\colon U\to\RR^3$ with $U=I\times(\RR\setdiff\{0\})$ given by
	\[\mb x(t,v) = \alpha(t) + v\alpha'(t).\]
	We have $\partial\mb x/\partial t = \alpha'(t)+v\alpha''(t)$
	and $\partial\mb x/\partial v = \alpha'(t)$, so
	\[\frac{\partial\mb x}{\partial t}\wedge\frac{\partial \mb x}{\partial v}
	= v\alpha''(t)\wedge \alpha'(t)\ne 0\]
	because the curvature of $\alpha$ at $t$ is $\kappa(t) =
	\frac{|\alpha'(t)\wedge\alpha''(t)|}{|\alpha'(t)|^3}\ne 0$.
	Thus $\mb x\colon U\to\RR^3$ is a regular parametrized surface
	whose trace consists of two connected pieces
	with common boundary $\alpha(I)$.
\end{example}

Observe also that a regular parametrized surface can intersect itself,
just like a regular curve can intersect itself
(exercise, find one such surface).

\begin{prop}
	Let $\mb x\colon U\subseteq\RR^2\to\RR^3$ be a regular parametrized surface.
	Then, for each $q\in U$, there is a neighborhood $V\subseteq\RR^2$ of $q$
	such that $\mb x(V)$ is a regular surface.
\end{prop}

\begin{exercise}
	\emph{Say the line, Bart!} (``Implicit function theorem.'')
\end{exercise}

\subsection{2.4 The tangent plane and the differential of a map}
Although we have defined what it means for a map to be differentiable,
we haven't actually defined a notion of derivative yet.

To start, we show that at every point of a regular surface $S$,
there is a well-defined tangent plane.
By a \vocab{tangent vector} to $S$ at a point $p$,
we mean the vector $\alpha'(0)$, where $\alpha\colon(-\eps,\eps)\to S$
is a differentiable parametrized curve such that $\alpha(0) = p$.

\begin{prop}
	Let $\mb x\colon U\subseteq\RR^2\to S$ be a parametrization
	of a neighborhood of a point $p$ in a regular surface $S$
	and let $q = \mb x\inv(p)$. Then the subspace
	\[d\mb x_q(\RR^2)\subseteq\RR^3\]
	coincides with the set of tangent vectors to $S$ at $p$.
\end{prop}

\begin{proof}
	On one hand, let $w$ be a tangent vector at $p$
	and let $\alpha\colon(-\eps,\eps)\to S$ be a corresponding
	differentiable parametrized curve with $\alpha(0)=p$, $\alpha'(0)=w$.
	Then $\beta = \mb x\inv\circ\alpha\colon(-\eps,\eps)\to U$
	is differentiable, so, by the chain rule,
	\[w = \alpha'(0) = (\mb x\circ \beta)'(0)
	= d\mb x_{\beta(0)}\cdot \beta'(0) = d\mb x_q(\beta'(0))
	\in d\mb x_q(\RR^2).\]
	
	On the other hand, let $w = d\mb x_q(v)$ for some $v\in\RR^2$.
	Then we obviously have $v = \gamma'(0)$ where
	$\gamma\colon (-\eps,\eps)\to U$ is given by $\gamma(t) = vt+q$,
	but then $w = \alpha'(0)$ where $\alpha = \mb x\circ\gamma$,
	so $w$ is a tangent vector (noting that $\alpha(0) = \mb x(\gamma(0))
	= \mb x(q) = p$).
\end{proof}

We will denote the tangent plane $d\mb x_q(\RR^2)$
to $S$ at $p$ by $T_p(S)$; note that the above proof implies that
$T_p(S)$ is invariant of the choice of $\mb x$, as the notation indicates.
The choice of parametrization determines a basis for $T_p(S)$,
specifically $\{\mb x_u(q), \mb x_v(q)\}$.
Then, if $w\in T_p(S)$ is the velocity vector $\alpha'(0)$
of a curve $\alpha = \mb x\circ\beta$ where $\beta\colon(-\eps,\eps)\to U$
is given by $\beta(t) = (u(t), v(t))$ and $\beta(0) = q = \mb x\inv(p)$, then
\[\alpha'(0) = (\mb x\circ \beta)'(0) = \mb x'(\beta(0))\beta'(0)
= \mb x_u(q)u'(0) + \mb x_v(q)v'(0) = w,\]
so the coordinates of $w$ in this basis are $(u'(0), v'(0))$.

Now we can talk bout the differential of a differentiable map between surfaces.
Let $S_1$ and $S_2$ be regular surfaces and
$\varphi\colon V\subseteq S_1\to S_2$ a differentiable map on
an open set $V$ in $S_1$. If $p\in V$, then $w\in T_p(S_1)$
is equal to $\alpha'(0)$ for some differentiable parametrized
$\alpha\colon(-\eps,\eps)\to V$ with $\alpha(0) = p$.
Then $\beta = \varphi\circ\alpha$ is a map with $\beta(0) = \varphi(p)$,
so $\beta'(0)$ is a vector in $T_{\varphi(p)}(S_2)$.

\begin{prop}
	Given $w$ in the situation above, the vector $\beta'(0)$
	is independent of $\alpha$.
	The map $d\varphi_p\colon T_p(S_1)\to T_{\varphi(p)}(S_2)$
	defined by $d\varphi_p(w) = \beta'(0)$ is linear.
\end{prop}

\begin{proof}
	Let $\mb x(u,v)$ and $\ol{\mb x}(\ol u,\ol v)$ be parametrizations
	in neighborhoods of $p$ and $\varphi(p)$ respectively.
	Then if $\varphi(u,v) = (\varphi_1(u,v), \varphi_2(u,v))$
	and $\alpha(t) = (u(t),v(t))$, we can calculate
	\[\beta'(0) = \left(\frac{\partial\varphi_1}{\partial u}u'(0)
	+ \frac{\partial\varphi_1}{\partial v}v'(0),
	\frac{\partial\varphi_2}{\partial u}u'(0)
	+ \frac{\partial\varphi_2}{\partial v}v'(0)\right),\]
	which are dependent only on $\varphi$ and the coordinates of $w$
	with respect to the basis induced by $\mb x$. Furthermore
	\[\beta'(0) = d\varphi_p(w) =\begin{pmatrix}
		\partial\varphi_1/\partial u&
		\partial\varphi_1/\partial v\\
		\partial\varphi_1/\partial u&
		\partial\varphi_1/\partial v
	\end{pmatrix}\begin{pmatrix}
		u'(0)\\v'(0)
	\end{pmatrix}\]
	shows that $d\varphi_p$ is clearly a linear map from
	$T_p(S_1)$ to $T_{\varphi(p)}(S_2)$.
\end{proof}

The map $d\varphi_p$ is the differential of $\varphi$ at $p\in S_1$.
We can similarly define the differential of a differentiable function
$f\colon U\subseteq S_1\to\RR$ at a point $p\in U$
as a linear map $df_p\colon T_p(S)\to\RR$.

We can also generalize the inverse function theorem.
\begin{definition}
	A mapping $\varphi\colon U\subseteq S_1\to S_2$ is
	a \vocab{local diffeomorphism} at $p\in U$
	if there exists a neighborhood $V\subseteq U$ of $p$
	such that $\varphi|_V$ is a diffeomorphism onto
	an open set $\varphi(V)\subseteq S_2$.
\end{definition}

\begin{prop}[Inverse function theorem]
	If $S_1$ and $S_2$ are regular surfaces and
	$\varphi\colon U\subseteq S_1\to S_2$ a differentiable map
	on an open set $U$ and $p\in U$ is a point
	such that $d\varphi_p$ is invertible,
	then $\varphi$ is a local diffeomorphism at $p$.
\end{prop}

We end this section by observing that
given parametrization $\mb x\colon U\subseteq\RR^2\to S$,
we can compute the normal vector at a given $q\in\mb x(U)$ by
\[N(q) = \frac{\mb x_u\wedge\mb x_v}{|\mb x_u\wedge\mb x_v|}(q).\]

\pagebreak
\subsection{2.5 The first fundamental form}
The inner product of $\RR^3$ induces an inner product
on each tangent space $T_p(S)$ in a natural way:
for $w_1,w_2\in T_p(S)$, $\langle w_1,w_2\rangle_p = w_1\cdot w_2$.
Then, we define the \vocab{first fundamental form} on $T_p(s)$
to be the quadratic form
\[I_p(w) = \langle w,w\rangle_p = |w|^2 > 0.\]
Given a parametrization $\mb x(u,v)$ at $p$,
we can express $I_p$ in terms of $\mb x$:
given $w\in T_p(S)$ and $\alpha$ a curve corresponding to $w$,
\begin{align*}
	I_p(\alpha'(0)) &= \langle\alpha'(0),\alpha'(0)\rangle_p\\
	&= \langle \mb x_uu'+\mb x_vv', \mb x_uu'+\mb x_vv'\rangle\\
	&= \langle\mb x_u,\mb x_u\rangle_p(u')^2
	+ 2\langle \mb x_u,\mb x_v\rangle_p u'v'
	+ \langle\mb x_v,\mb x_v\rangle_p(v')^2\\
	&= E(u')^2 + 2Fu'v' + G(v')^2
\end{align*}
where
\[E(u_0,v_0):=\langle \mb x_u,\mb x_u\rangle_p,\quad
F(u_0,v_0):=\langle \mb x_u,\mb x_v\rangle_p,\quad
G(u_0,v_0):=\langle \mb x_v,\mb x_v\rangle_p.\]

\begin{example}
	If $S$ is a plane passing through $p_0$ and containing orthonormal vectors
	$w_1$, $w_2$, then the plane can be parametrized
	\[\mb x(u,v) = p_0 + uw_1 + vw_2,\]
	so we can compute $(E,F,G) = (1,0,1)$ everywhere.
\end{example}

\begin{example}
	Given the right cylinder $x^2+y^2=1$, we have the parametrization
	$\mb x\colon(0,2\pi)\times\RR\to\RR^3$ everywhere except for the line
	$(1,0,0)+(0,0,1)t$ given by
	\[\mb x(u,v) = (\cos u,\sin u,v),\]
	so we can compute $(E,F,G) = 1$ again.
\end{example}

\begin{exercise}
	Compute $E$, $F$, and $G$ for the parametrization of the helicoid
	given by $(v\cos u,v\sin u,au)$ for $0<v<2\pi$, $-\infty<v<\infty$.
\end{exercise}

The first fundamental form is so important because, abstractly,
knowing how $I$ behaves lets you understand the metric properties of $S$
without talking about the ambient space $\RR^3$.
Of course, it's still important that we can compute things
by using a parametrization and writing expressions in $E$, $F$, $G$.

\begin{exercise}
	Let $\alpha(t) = \mb x(u(t),v(t))$ be a curve contained in
	a coordinate neighborhood corresponding to the parametrization $\mb x$.
	Show that the arc length of $\alpha$ between $0$ and $t$ is
	\[s(t) = \int_0^t\sqrt{E(u')^2 + 2Fu'v' + G(v')^2}\dif x.\]
\end{exercise}

\begin{exercise}
	Show that the angle $\varphi$ between the coordinate curves of
	a parametrization $\mb x\colon U\to S$ at $q\in U$ satisfies
	\[\cos\varphi = \frac{\langle\mb x_u,\mb x_v\rangle}{|\mb x_u||\mb x_v|}(q)
	= \frac{F}{\sqrt{EG}}.\]
\end{exercise}

We can also deal with area: for a compact subset of $\RR^2$ contained
in a coordinate neighborhood $\mb x\colon U\to S$
(so $R := \mb x(Q)$ is a bounded region in $S$),
the quantity
\[\iint_Q |\mb x_u\wedge\mb x_v|\dif u\dif v\]
can be shown to be independent of $\mb x$ via the change of variables formula.
We thus take as a definition that the area of $R$ is
\[A(R) := \iint_Q |\mb x_u\wedge\mb x_v|\dif u\dif v.\]
\begin{exercise}
	Show that $|\mb x_u\wedge\mb x_v| = \sqrt{EG - F^2}$.
	Then, parametrizing a torus as
	\[\mb x(u,v) = ((a+r\cos u)\cos v,(a+r\cos u)\sin v,r\sin u)\]
	for $(u,v)\in(0,2\pi)\times(0,2\pi)$ everywhere except for
	a single meridian and a single parallel, show that
	the surface area of the torus is $4\pi^2 ra$.
	Here $r$ is the radius on which the vertical cross-sectional
	circles' centers vary, and $a$ is the radius of
	the vertical cross-sectional circles.
\end{exercise}

\begin{exercise}
	Parametrize a sphere with spherical coordinates and then
	show that a sphere's surface area is $4\pi r^2$.
\end{exercise}

\section{The Geometry of the Gauss Map}
\subsection{3.2 The Gauss Map}
Given a coordinate neighborhood $\mb x(U)$, recall that
the unit normal vector $N\colon \mb x(U)\to S^2$ at a point $q$ is given by
the differentiable function
\[N(q) = \frac{\mb x_u\wedge \mb x_v}{|\mb x_u\wedge\mb x_v|}.\]
Note that this may not necessarily be able to be extended to the entire surface
(consider for example the M\"obius strip).
If the map is able to be extended to all of $S$, then we say that
$S$ is \vocab{orientable}, and we call the vector field $N$
an \vocab{orientation} of $S$.
An orientation on $S$ then induces an orientation on the tangent space
at some point $p\in S$ in that a basis $\{v,w\}\subseteq T_p(S)$
can be called \vocab{positive} if $\langle v\wedge w,N\rangle$ is positive.

From this point on let $S$ be a regular orientable surface
with an orientation $N$; we will frequently omit
the adjectives ``regular orientable'' for brevity.

\begin{definition}[Gauss map]
	Let $S$ be a surface with orientation $N$.
	The map $N\colon S\to S^2$ is called the \vocab{Gauss map} of $S$.
\end{definition}

We can easily verify that the Gauss map is differentiable,
and its differential is a linear map from $T_p(S)$ to $T_{N(p)}(S^2)$,
but these two vector spaces are the same (since the normal to $S$ at $p$
is also a normal to $S^2$ at $N(p)$), so we can think of $dN_p$
as an endomorphism of $T_p(S)$.
Moreover, the differential measures the rate
at which $N$ is pulling away from $N(p)$ in a neighborhood of $p$ ---
this is the analog of the curvature for surfaces.

\begin{prop}
	The differential $dN_p\colon T_p(S)\to T_p(S)$ of the Gauss map
	is a self-adjoint linear map.
\end{prop}

\begin{proof}
	To show this we need to show that for any $w_1,w_2\in T_p(S)$,
	\[\langle dN_p(w_1),w_2\rangle = \langle w_1,dN_p(w_2)\rangle\]
	so in particular it suffices to show this for some basis $w_1$, $w_2$.
	If $\mb x$ is a parametrization around $p$, then this reduces to showing
	\[\langle N_u,\mb x_v\rangle = \langle \mb x_u, N_v\rangle.\]
	This follows from e.g. writing down the derivative of
	$\langle N,\mb x_u\rangle = 0$ with respect to $v$ and the derivative of
	$\langle N,\mb x_v\rangle = 0$ with respect to $u$.
\end{proof}

\begin{definition}
	The \vocab{second fundamental form} of $S$ at $p$ is
	the quadratic form $II_p(v)$ on $T_p(S)$ defined as
	\[II_p(v) = -\langle dN_p(v),v\rangle.\]
\end{definition}





















\end{document}